% 01_contexto.tex - Contexto, estado del arte y oportunidad.
\section{Contexto, estado del arte y oportunidad}
% TODO: encuadre del problema y el sector; estado del arte con referencias
% verificadas; y oportunidad concreta que abre la colaboracion
% publico-privada.

\subsection{Contexto y problema}
% TODO: descripcion del problema tecnico o social, quien lo sufre y por que
% las soluciones actuales no lo resuelven. Estilo sobrio y directo
% (regla anti-fluff del marco academico).

\subsection{Estado del arte}
% TODO: sintesis del estado del arte apoyada en:
% - Las 4 fuentes MCP (Semantic Scholar, arXiv, PubMed, CrossRef), registradas
%   por lit-synthesizer en src/references.bib (ver seccion 6 del marco).
% - El conocimiento local del consorcio en rag/past_projects/ (inspiracion
%   estructural; los datos viejos se REVERIFICAN contra el SSOT antes de
%   usarse, ver seccion 8 del marco).
% Citar SOLO claves existentes en src/references.bib: si falta evidencia,
% \todo{CITA REQUERIDA: [que evidencia se necesita]}.

\subsection{Oportunidad y encaje en la convocatoria}
% TODO: por que este consorcio y por que ahora; encaje con los objetivos y
% prioridades de la convocatoria (requisitos y pesos en
% context/convocatoria_checklist.md y context/call_guidelines.md).
